\section{The unchanged physical problem}
Put $R_-=9/20$ and $R_+=47/100$, and let $R\in I=[R_-,R_+]$. Circular scatterers have centers in
$\Lambda=\Z(1,0)+\Z(1/2,\sqrt3/2)$. Write $T_R$ for the regular
collision map and $(\alpha,p)\in(\R/2\pi\Z)\times(-1,1)$ for its
coordinates; $p$ is the outgoing tangential velocity and
$\dd\nu=\dd\alpha\,\dd p/(4\pi)$. Physical flights satisfy
$\tau_-\le\tau_R\le\tau_+$, where $\tau_-=3/50$ and $\tau_+$
is the fixed finite-horizon bound.

The section is specified here to fix every event in the argument.
Put $h_0=1/200$. Let $Y_R$ be the union of the four open coordinate
squares of half-side $h_0$ centered at
\[
 (0,0),\quad(\pi/6,-1/2),\quad
 (-\theta_R,\sin\theta_R),\quad
 (\pi/3-\theta_R,\sin\theta_R).
\]
Here $\theta_R=2\arctan q_R$, and $q_R$ is the unique root in
$(49/100,53/100)$ of
\[
 8Rq(1+q^2)-4q(1-q^2)+\sqrt3(1-6q^2+q^4)=0.
\]
Set
\[
 U=\{|\alpha-\pi/6|<3/50,\ |p|<3/20\},\qquad
 Y_R^*=Y_R\cup U,\qquad c_*=\frac{91}{10000\pi}.
\]
The rectangles are disjoint and $\nu(Y_R^*)=c_*$. We always use
$\nu_R^*=c_*^{-1}\nu|_{Y_R^*}$, not a renormalization of a smaller
source.

At $m$ flights, let $L=L_{m,R}$ be their total length,
$K^{\rm c}_{m,R}$ their lattice displacement, and
$A_{m,R}=\sum_{j=0}^{m-1}\one_{Y_R^*}\circ T_R^j$. The event is
\begin{equation}\label{eq:v76-event}
 \Pi_{n,k,m,R}=Y_R^*\cap\{K^{\rm c}_{m,R}=k,\ A_{m,R}=n\}
                         \cap T_R^{-m}Y_R^*.
\end{equation}
The initial visit is counted; the terminal visit is tested separately.
On this event the actual $n$th return has collision count $m$ and
roof $L$. Denote its density by $p_{n,R}(k,m,t)$.

Write $c_j=\sqrt{1-p_j^2}$ and let $\Delta_j$ be the least clearance
of flight $j$ from any third disk, excluding its two incident disks.
Keep the original smooth function $0\le\Theta\le1$, zero on
$(-\infty,1]$ and one on $[2,\infty)$. With $q_0=47/53$,
$d_j=\min(j,m-j)$ and $d'_j=\min(j,m-1-j)$, set
\begin{align}
 H^\varepsilon&=
 \prod_{j=0}^{m}\Theta\!\left(\frac{c_j}{\varepsilon q_0^{d_j/4}}\right)
 \prod_{j=0}^{m-1}\Theta\!\left(\frac{\Delta_j}{\varepsilon q_0^{d'_j/4}}\right),
 \qquad W^\varepsilon=1-H^\varepsilon,\label{eq:v76-guard}\\
 b^\varepsilon(t)\dd t
   &=L_\#(W^\varepsilon\one_{\Pi_{n,k,m,R}}\nu_R^*).
 \label{eq:v76-source}
\end{align}
In particular first incidence and first clearance outside any selected
critical neighborhoods are still present. No section-distance factor
has been inserted into this guard.

\subsection{The arithmetic reference and the target}
For $y=(k_1,k_2,t,n)$ put
\[
 \mu_R=(0,0,\bar\tau_R,c_*),\quad
 Z=(y-m\mu_R)/\sqrt m,\quad
 \mathcal C_{m,R,n,k}(M)=\{t:|Z|\le M\}.
\]
We retain the finite spectral-peak reference, including peaks close to
but not on the unit circle. To specify its meaning, the collision
record observable is $(\kappa_1,\kappa_2,\tau,\one_{Y_R^*})$.
Let $\lambda_j(R,\vartheta)$ and $\mathsf P_j(R,\vartheta)$ be its
continued simple twisted eigenvalue and projection in each of the
fixed peak charts, and let $\chi_j$ be their frequency partition.
At the real maximum $a_j$ of $\log|\lambda_j|$ put
\[
 z_j=\lambda_j(R,a_j),\quad
 \mu_j=i^{-1}\nabla\log\lambda_j(R,a_j),\quad
 H_j=-\nabla^2\log\lambda_j(R,a_j),
\]
with $\operatorname{Re}H_j\ge aI_4$ for a common $a>0$. If
$\eta_R=\one_{Y_R^*}$ and $\ell$ denotes mass evaluation, put
$B_j=\chi_j(R,a_j)\ell(M_{\eta_R}\mathsf P_j(R,a_j)(\eta_R\nu))$.
Define
\begin{align}
 G_H(z)&=(2\pi)^{-4}\int_{\R^4}e^{-iv\cdot z-v^{\mathsf T}Hv/2}\dd v,
 \notag\\
 \mathcal L_{m,R}(y)&=\operatorname{Re}\left[
 c_*^{-1}\sum_j e^{-ia_j\cdot y}z_j^m B_j
      G_{H_j}\left((y-m\mu_j)/\sqrt m\right)\right].
 \label{eq:v76-reference}
\end{align}
Only the charts active at $R$ are summed. This is the same transition
reference as in the retained spectral theorem. The existence and
uniform estimates for these scalar data belong to that analytic
input, not to the transport argument below. No frequency branch or
arithmetic zero class is deleted.

The governing assertion remains
\begin{equation}\label{eq:v76-target}
 \sup_{R,\,1\le n\le m,\,k\in\Z^2}
 \esssup_{t\in\mathcal C_{m,R,n,k}(M)}
 |m^2p_{n,R}(k,m,t)-\mathcal L_{m,R}(k_1,k_2,t,n)|\longrightarrow0.
\end{equation}
Let $\mathcal E_M$ denote this error's collision limsup. An empty
central target contributes zero. The order of the uniformity is
$\sup_R\esssup_t$, so exceptional density sets may depend on $R$.


